\section{ТЕОРЕТИЧНІ ЗАСАДИ НАВЧАННЯ РОЗРОБКИ ІМЕРСИВНИХ ОСВІТНІХ РЕСУРСІВ}

У розділі висвітлено концептуальні основи навчання розробки імерсивних освітніх ресурсів майбутніх учителів інформатики, що постають одним із ключових напрямів педагогічної інноватики в умовах цифрової трансформації освіти. Розкрито педагогічну проблематику створення VR/AR/XR-рішень, що потребує від майбутнього вчителя інтеграції технологічних, методичних та дизайнерських компетентностей у межах цілісної професійної картини діяльності. Уточнено сутнісні характеристики імерсивних освітніх ресурсів, які представлені як інтерактивні засоби нового покоління, побудовані на моделюванні віртуальних або доповнених середовищ та спрямовані на підсилення дидактичних можливостей навчання. Здійснено систематизацію вітчизняного досвіду використання імерсивних технологій в освітньому процесі, зосереджену на аналізі освітньо-професійних програм підготовки майбутніх учителів інформатики.

\subsection{Розробка імерсивних освітніх ресурсів як педагогічна проблема підготовки майбутніх учителів інформатики}\label{subsec:1.1}

Активне впровадження новітніх педагогічних технологій, які змінюють характер взаємодії учасників освітнього процесу, - провідна ознака часу. Освіта XXI ст. розвивається під впливом процесів цифровізації, швидкого оновлення знань і зростання ролі технологій, які забезпечують персоналізоване, адаптивне та інтерактивне навчання \cite{UNESCO2021DigEd}.

Посилення ролі ІКТ закономірно спрямовує увагу на імерсивні технології як засоби поглибленого занурення у навчальний зміст і розширення можливостей професійної підготовки фахівців. За таких умов вивчення імерсивних технологій набуває особливої актуальності у підготовці майбутніх вчителів, зокрема вчителів інформатики, оскільки ІТ формують новий рівень взаємодії з навчальним середовищем, забезпечуючи занурення у зміст і розширюючи можливості педагогічної діяльності в умовах цифрової трансформації освіти. Так імерсивні освітні ресурси (далі - ІОР) стають невід'ємним складником цифрової трансформації освіти, що відбувається в умовах глобального інформаційного суспільства і формує нові вимоги до професійної діяльності педагога. Цифрова компетентність сьогодні є необхідною умовою професійної діяльності сучасного вчителя, оскільки вона є запорукою ефективного виконання вчителем педагогічних завдань та досягнення освітніх цілей~\cite[c.~194]{Kibenko2022Digital}. Отже проблема породжує потребу у підготовці майбутніх учителів до формування цифрової компетентності, а особливо використання імерсивних технологій, апробацію нововведень та їх реалізацію в умовах сучасних закладів загальної середньої освіти~\cite{Kibenko2022Digital}.

За таких умов учитель, зокрема вчитель інформатики, має бути не лише користувачем готових цифрових рішень, а й творцем інноваційних освітніх середовищ, здатним інтегрувати технології доповненої, віртуальної та змішаної реальності у власну педагогічну практику. Результати систематичного огляду та мета-аналізу Х.~Гарзона, Дж. Павона та С.~Балдіріса засвідчують статистично значущий позитивний вплив доповненої реальності на навчальні результати в різних освітніх контекстах \cite{Akdeniz2023AR}. У підготовці майбутніх учителів інформатики розробка таких ресурсів виходить за межі суто технічного завдання, набуваючи статусу психолого-педагогічної проблеми, оскільки передбачає формування у студентів не лише цифрових навичок, а й педагогічного мислення, креативності, рефлексивності та етичної відповідальності за якість освітнього продукту \cite{Kerr2020Augmented}.

З позицій сучасної педагогічної науки імерсивне навчання розглядається як діяльність, що створює умови для розвитку вищих психічних функцій через досвід занурення у віртуально-реальні ситуації, наближені до реальних умов майбутньої професійної діяльності. Психолого-педагогічний аспект розробки ІОР зумовлений необхідністю врахування закономірностей сприйняття, уваги, пам'яті, мотивації та емоційного залучення користувача. Від рівня усвідомлення цих процесів залежить ефективність побудови навчального сценарію, дизайну середовища та інтеграції інструментів VR/AR/MR у педагогічну практику \cite{Radianti2020VR}.

Проблемність теми полягає у відсутності усталених методик формування готовності майбутніх учителів до створення ІОР, невизначеності критеріїв їхньої педагогічної ефективності та недостатній інтеграції психолого-педагогічних підходів до цифрового дизайну.

На цьому тлі актуалізується пошук педагогічних умов і методичних стратегій, що сприятимуть розвитку в майбутніх учителів інформатики здатності до свідомого, творчого й етично виваженого використання технологій занурення в освітньому процесі~\cite{Goncharenko2011Pedagogy}. Дослідники зазначають, що у процесі підготовки майбутніх учителів інформатики розробка імерсивних освітніх ресурсів виступає одночасно засобом і результатом формування професійних компетентностей \cite{Alalwan2020VR}. За таких умов студент постає не лише користувачем цифрових технологій, а проєктувальником інтерактивного освітнього середовища, здатним забезпечити педагогічну доцільність, доступність і безпечність цифрового контенту. Отже, навчання розробки імерсивних освітніх ресурсів потребує цілісного психолого-педагогічного підґрунтя, що інтегрує положення когнітивної психології, педагогічного дизайну, дидактики та теорії навчання інформатики \cite{Dalgarno2010What}.

%%% PART A: Theoretical Problem Framing %%%

Теоретичне осмислення проблеми навчання розробки ІОР потребує опори на провідні міжнародні моделі, що розкривають механізми навчання в імерсивних середовищах. Когнітивно-афективна модель імерсивного навчання (CAMIL), запропонована Г.~Макранскі та Г.~Б.~Петерсеном~\cite{Makransky2021CAMIL}, визначає шість ключових факторів, що опосередковують вплив імерсивних технологій на навчальні результати: присутність (presence), агентність (agency), втілення (embodiment), когнітивне навантаження, мотивацію та саморегуляцію. Модель CAMIL обґрунтовує, що ефективність імерсивного навчання зумовлюється не самою технологією, а дизайном освітнього досвіду, який має враховувати когнітивні та афективні процеси учасників. Це положення є концептуальною передумовою для обґрунтування необхідності спеціальної підготовки вчителів, оскільки педагогічно виважене проєктування ІОР потребує розуміння зазначених психолого-когнітивних механізмів.

Теоретичну основу проєктування ефективних імерсивних ресурсів доповнює теорія когнітивного навантаження. У систематичному огляді Й.~Бухнера, К.~Бунтінс та М.~Керреса~\cite{Buchner2021CogLoad} встановлено, що доповнена реальність здатна знижувати когнітивне навантаження порівняно з традиційними мультимедійними засобами, проте цей ефект суттєво залежить від дизайну навчальних матеріалів. М.~Пупар зі співавторами~\cite{Poupard2024Systematic} на основі систематичного огляду 36 досліджень виявили діаметрально протилежні тенденції: віртуальна реальність збільшує стороннє когнітивне навантаження, тоді як доповнена реальність оптимізує його, що робить AR більш придатною для навчання початківців. Це має пряме значення для підготовки майбутніх учителів, які зазвичай є початківцями у роботі з імерсивними технологіями. Я.~Гоннерманн-Мюллер та Ю.~М.~Крюгер~\cite{GonnermannMuller2024ARDesign} систематизували 15 доказових рекомендацій щодо дизайну AR-ресурсів за шістьма вимірами, підкреслюючи, що розробник ІОР має оперувати не лише технічними засобами, а й принципами управління когнітивним навантаженням.

Водночас аналіз досліджень виявляє глибоку суперечність між визнанням ефективності імерсивних технологій та станом підготовки педагогів до їх створення. Х.~Бельда-Медіна та Х.~Р.~Кальво-Феррер~\cite{BeldaMedina2022TPACK} у дослідженні інтеграції доповненої реальності в мовну освіту крізь призму моделі TPACK встановили, що майбутні вчителі демонструють критичний дефіцит на перетині технологічних і педагогічних знань (TPK), оскільки їх готують переважно як \emph{користувачів} AR-засобів, а не як їх \emph{розробників}. С.~Обердорфер зі співавторами~\cite{Oberdorfer2021Interdisciplinary} експериментально підтвердили, що міждисциплінарна модель навчання, яка поєднує студентів педагогічних та ІТ-спеціальностей, забезпечує значне покращення як медіадидактичних, так і технологічних знань обох груп, засвідчуючи потенціал інтеграційного підходу до підготовки вчителів-розробників ІОР.

Ефективність впровадження імерсивних технологій в освіту значною мірою визначається готовністю педагогів до їх використання та створення. Дж.~Янг зі співавторами~\cite{Jang2021TAM} на основі розширеної моделі прийняття технології (TAM) та TPACK встановили, що рівень технологічних та педагогічних знань учителя статистично значуще впливає на сприйняту корисність і намір використовувати AR/VR-технології. Л.~Тан зі співавторами~\cite{Tan2023TPACK} за результатами кластерного аналізу TPACK австралійських учителів виявили стійкі компетентнісні розриви у сфері цифрової грамотності, зокрема у навичках створення мультимодальних ресурсів, що підтверджує актуальність цілеспрямованої підготовки вчителів до розробки ІОР.

Узагальнення результатів провідних мета-аналізів (табл.~\ref{tab:meta_analysis}) засвідчує, що імерсивні технології мають стійкий статистично значущий позитивний вплив на навчальні результати (розміри ефекту від g\,=\,0{,}82 до d\,=\,0{,}98) у різних освітніх контекстах і галузях.

\begin{table}[!h]
\centering
\caption{Узагальнення результатів мета-аналізів ефективності імерсивних технологій в освіті}
\label{tab:meta_analysis}
\small
\begin{tabular}{|l|c|l|l|l|}
\hline
\hfill \textbf{Мета-аналіз} & \textbf{\textit{N}} & \textbf{Технологія} & \hfill \textbf{Розмір ефекту} & \hfill \textbf{Галузь} \\
\hline
Zhang et al. \cite{Omelchuk2024Secondary} & 129 & AR, K-12 & g\,=\,0{,}919 & STEM \\
\hline
Coban et al. \cite{Lytvynova2023Comparative} & 48 & IVR & d\,=\,0{,}98 & Вища освіта \\
\hline
Chen et al. \cite{Boiko2024Immersive} & 17 & XR & g\,=\,0{,}825 & Мовна освіта \\
\hline
Kindzhibala et al. \cite{Kindzhibala2025Creative} & 19 & VR/AR & g\,=\,0{,}82 & Креативність \\
\hline
Buchner et al. \cite{Buchner2021CogLoad} & 58 & AR & зниження КН$^{*}$ & Загальний \\
\hline
Poupard et al. \cite{Poupard2024Systematic} & 36 & VR/AR & VR↑ КН, AR↓ КН & Загальний \\
\hline
\multicolumn{5}{l}{\footnotesize $^{*}$КН~-- когнітивне навантаження} \\
\end{tabular}
\end{table}

%%% PART B: Bibliometric Dynamics %%%

Таким чином, провідні міжнародні моделі імерсивного навчання (CAMIL, теорія когнітивного навантаження, TPACK) та результати мета-аналізів засвідчують теоретичну обґрунтованість і практичну затребуваність проблеми підготовки вчителів до розробки ІОР. Актуальність цієї проблеми підтверджується також динамікою наукових публікацій, що відображає стрімке зростання дослідницького інтересу до імерсивних технологій в освіті.

Динаміку наукового інтересу засвідчує бібліометричний аналіз. С.~Гарнал зі співавторами~\cite{Harnal2023Bibliometric} здійснили бібліометричне картування 3909 публікацій з проблематики доповненої реальності в освіті, виявивши стійке зростання дослідницької активності з чітко вираженим прискоренням після 2018 року та визначивши чотири основні тематичні напрями: дидактичні підходи, технологічні рішення, педагогічну підготовку та інклюзивне застосування. Аналіз розподілу джерельної бази дисертації за роками публікації (рис.~\ref{fig:biblio_dynamics}) підтверджує цю тенденцію: переважна більшість використаних джерел припадає на період 2019--2023 років, що відображає стрімке зростання наукового інтересу до імерсивних технологій в освіті та актуальність порушеної проблематики.

\begin{figure}[!h]
\centering
\begin{tikzpicture}
\begin{axis}[
    compat=newest,
    axis lines=left,
    ybar,
    bar width=14pt,
    enlargelimits=0.05,
    width=\textwidth,
    height=0.45\textwidth,
    xlabel={\textbf{рік}},
    ylabel={\textbf{кількість публікацій}},
    symbolic x coords={2001, 03, 09, 2010, 11, 12, 13, 14, 15, 16, 17, 18, 19, 2020, 21, 22, 23},
    xtick=data,
    nodes near coords,
    every node near coord/.append style={font=\scriptsize},
    ymajorgrids=true,
    grid style={dashed, gray!30},
]
\addplot[fill=black!60] coordinates {
    (2001,1) (03,2) (09,2) (2010,2) (11,1) (12,1) (13,1) (14,5) (15,6)
    (16,4) (17,2) (18,6) (19,13)
    (2020,16) (21,16) (22,14) (23,1)
};
\end{axis}
\end{tikzpicture}

\caption{Розподіл джерельної бази дисертації за роками публікації}
\label{fig:biblio_dynamics}
\end{figure}

%%% Cluster Analysis Introduction %%%

Аналіз джерельної бази засвідчує доцільність її систематизації за кількома укрупненими проблемно-тематичними кластерами, що відображають основні напрями науково-педагогічного осмислення імерсивних технологій в освіті (рис.~\ref{fig:1.1}).

\begin{figure}[!h]
\centering
%\resizebox{\textwidth}{!}{%
\begin{tikzpicture}

% Title
\node[dia/header=14cm] (title) {ПРОБЛЕМНО-ТЕМАТИЧНІ КЛАСТЕРИ};

% 2x2 grid of cluster boxes
\node[dia/rbox=5.5cm, below=1.5cm of title.south, xshift=-4cm] (c1)
  {Загальнотеоретичний і дидактичний кластер};
\node[dia/rbox=5.5cm, below=1.5cm of title.south, xshift=4cm] (c2)
  {Освітні практики};
\node[dia/rbox=5.5cm, below=1.2cm of c1] (c3)
  {Галузеві, порівняльні та узагальнювальні підходи};
\node[dia/rbox=5.5cm, below=1.2cm of c2] (c4)
  {Професійна педагогічна підготовка};

% Tree structure: title → row 1, row 1 → row 2
\draw[dia/line] (title.south) -- (c1.north);
\draw[dia/line] (title.south) -- (c2.north);
\draw[dia/line] (c1.south) -- (c3.north);
\draw[dia/line] (c2.south) -- (c4.north);

% Convergence point
\coordinate (conv) at ($(c3.south)!0.5!(c4.south) + (0,-1.5cm)$);
\node[dia/label] at (conv) (convlabel) {Точка перетину};

% Lines from row 2 boxes to convergence (direct, no overlap)
\draw[dia/line] (c3.south) -- (convlabel.north);
\draw[dia/line] (c4.south) -- (convlabel.north);

% Lines from row 1 boxes to convergence (routed outside row 2)
\draw[dia/line] (c1.south west) -- ([xshift=-0.4cm]c3.south west) -- (convlabel.north west);
\draw[dia/line] (c2.south east) -- ([xshift=0.4cm]c4.south east) -- (convlabel.north east);

% Bottom box
\node[dia/rbox=6cm, below=1.2cm of convlabel] (inclusive)
  {Інклюзивний освітній вимір};

% Arrow from convergence to bottom box
\draw[dia/arrow] (convlabel.south) -- (inclusive.north);

\end{tikzpicture}%
%}

\caption{Проблемно-тематичні кластери науково-педагогічного осмислення імерсивних технологій в освіті}
\label{fig:1.1}
\end{figure}

Перший кластер охоплює \emph{{загальнотеоретичні та дидактичні підходи до трактування імерсивних освітніх ресурсів}.} У межах цього напряму імерсивні технології розглядаються як інноваційний дидактичний інструмент, що зумовлює трансформацію змісту, форм і методів навчання, розширює можливості організації освітнього процесу та сприяє реалізації сучасних моделей навчання, орієнтованих на гнучкість і варіативність освітнього середовища.

Другий кластер репрезентує \emph{{практики впровадження імерсивних технологій у освітній процес на різних рівнях і в різних ланках освіти}.} У цьому контексті увага зосереджується на педагогічних умовах, організаційних моделях і дидактичних сценаріях використання імерсивних ресурсів з урахуванням вікових, когнітивних і освітніх особливостей здобувачів освіти.

Третій блок об'єднує \emph{\textbf{дослідження використання імерсивних технологій у системі професійної педагогічної підготовки та розвитку педагогічного професіоналізму}.} Імерсивні ресурси осмислюються як засіб формування професійних і цифрових компетентностей майбутніх учителів, розвитку їхньої методичної готовності та здатності до проєктування освітніх рішень у цифровому освітньому середовищі.

Четвертий кластер становлять \emph{{дослідження застосування імерсивних технологій в інклюзивній освіті}.} У межах цього напряму імерсивні освітні ресурси розглядаються як інструмент забезпечення доступності навчання, підтримки різноманітності освітніх потреб та створення безбар'єрного, адаптивного освітнього середовища для всіх категорій здобувачів освіти.

П'ятий кластер охоплює \emph{{галузеві, порівняльні та узагальнювальні підходи до використання імерсивних технологій в освіті}.} У цьому блоці імерсивні ресурси аналізуються як дидактичний засіб розвитку предметних і професійних компетентностей у різних освітніх галузях, а також у контексті порівняльного аналізу міжнародного й вітчизняного досвіду та узагальнення результатів наукових і прикладних досліджень.

Коротко окреслимо коло аспектів, які набули найбільшої уваги у працях вітчизняних і зарубіжних науковців та визначають провідні напрями осмислення означеної проблеми у сучасному науково-педагогічному дискурсі.

%%% CLUSTER 1 %%%

У межах \emph{{першого кластера наукових джерел}} актуалізується проблематика класифікації імерсійних технологій, визначення їх дидактичного потенціалу та окреслення міждисциплінарних можливостей застосування в освітньому процесі. Дослідники зосереджуються на осмисленні VR/AR-технологій як інноваційного дидактичного інструментарію, здатного трансформувати зміст навчання, розширити спектр освітніх сценаріїв і підсилити пізнавальну активність здобувачів освіти.

У контексті розглянутих вище теоретичних засад праці першого кластера збагачуються дослідженнями механізмів втілення та переносу навчання в імерсивних середовищах. С.~Клінгенберг зі співавторами~\cite{Klingenberg2024Embodiment} на основі моделі «занурення--інтерактивність» експериментально встановили, що втілення у віртуальній реальності сприяє переносу навчання, однак цей ефект модерується рівнем інтерактивності та структурою навчального завдання, що підтверджує положення CAMIL~\cite{Makransky2021CAMIL} щодо визначальної ролі педагогічного дизайну імерсивного досвіду.

Значний внесок у розроблення класифікаційних підходів до імерсивних технологій здійснено у працях Ю.~Носенко \cite{Nosenko2024Typology}, де запропоновано структурну типологію імерсивних сервісів для освітнього процесу. Авторка систематизує занурювальні середовища за низкою критеріїв, серед яких провідними виступають педагогічна мета використання, рівень інтерактивності, спосіб візуалізації та технічна складність. У роботі наголошується, що результативність упровадження імерсивних технологій визначається не лише технічними характеристиками, а й узгодженістю цифрових інструментів з дидактичними завданнями та готовністю користувачів до взаємодії з віртуальними моделями. Проблема \emph{персоналізації навчання в умовах цифрового середовища} розкривається у систематичному огляді А.~Маругкаса зі співавторами \cite{Gurska2025Personalized}, який охоплює 69 досліджень і присвячений питанню ефективності та персоналізації імерсивного VR-навчання. Автори систематизують механізми адаптації віртуальних освітніх середовищ до індивідуальних потреб здобувачів освіти, акцентуючи увагу на можливостях автоматизованого аналізу навчального прогресу та адаптації контенту відповідно до індивідуального темпу засвоєння матеріалу.

Окремий напрям у межах цього кластера становлять праці, присвячені \emph{дидактичному потенціалу доповненої реальності у шкільній освіті.} Зокрема, В.~Коваленко, М.~Мар'єнко та А.~Сухіх \cite{Pinchuk2022AR} доводять, що AR-контент забезпечує достовірне моделювання природничих процесів, поєднуючи візуалізацію з дослідницькою діяльністю учнів. Автори підкреслюють можливість використання доповненої реальності як засобу активізації навчального пізнання та формування дослідницьких умінь, наводячи приклади цифрових ресурсів, адаптованих до потреб шкільного навчання.

Розвиток ідей синтетичного та гібридного освітнього середовища послідовно представлено у працях Ю.~Богачкова \cite{Bogachkov2023Immersive} Автор вводить поняття «імерсивний синтетичний простір», у межах якого поєднуються очні й дистанційні форми навчання. У дослідженнях підкреслюється, що такі середовища сприяють глибшому зануренню в навчальну діяльність, забезпечують гнучкість освітніх сценаріїв і розширюють можливості індивідуального та групового навчання.

Психолого-педагогічний вимір застосування імерсивних технологій розкрито в експериментальному дослідженні Й.~Ванга та Ц.~Пань \cite{Lytvynova2025Emotional}, присвяченому впливу імерсивної віртуальної реальності на емоційний інтелект студентів університетів. Згідно з отриманими результатами, систематичне використання VR-середовищ сприяє розвитку здатності до усвідомлення, регуляції та конструктивного вираження емоцій, що відкриває перспективи формування психологічно орієнтованої цифрової педагогіки.

У контексті проєктування електронних освітніх ресурсів А.~Крассман зі співавторами \cite{Chornous2025Presence} на основі систематичного огляду 140 досліджень за 24 роки аналізують взаємозв'язок між відчуттям присутності та навчальними результатами у VR-середовищах. Автори встановлюють, що імерсивні технології здатні посилювати ефект присутності, покращувати запам'ятовування навчального матеріалу та формувати відчуття «дії від першої особи», наголошуючи на доцільності інтеграції VR-практик у систему цифрової дидактики. Подібну проблематику порушує й О.~Орлов \cite{Orlov2024Innovative} який розглядає імерсивні технології як елемент інноваційних освітніх маршрутів, що поєднують інтерактивність, візуальну насиченість і варіативність навчання. Автор акцентує на необхідності методичної підтримки педагогів для ефективного використання VR/AR-інструментів у практиці навчання.

Праці першого кластера формують цілісне уявлення про імерсивні технології як багатовимірний дидактичний феномен. У поєднанні з моделлю CAMIL~\cite{Makransky2021CAMIL} та концепцією «занурення--інтерактивність»~\cite{Klingenberg2024Embodiment} ці дослідження створюють теоретичне підґрунтя для аналізу практичної реалізації імерсивних технологій в освітньому процесі, який представлено у працях другого кластера.

%%% CLUSTER 2 %%%

У межах \emph{другого кластера} наукових джерел у центрі уваги перебуває проблематика практичного впровадження імерсивних технологій в освітній процес на різних рівнях і в різних ланках освіти, а також визначення педагогічних, методичних і технічних умов їх ефективного використання. Дослідження цього напряму репрезентують перехід від теоретичного осмислення імерсивних ресурсів до аналізу реальних освітніх практик, що охоплюють початкову, базову, профільну середню та вищу освіту, а також різні освітні галузі. При цьому звертається увага не лише на позитивні результати використання ІОР, а й на проблемні аспекти.

Систематичний огляд і мета-аналіз Ж.~Чжана зі співавторами \cite{Omelchuk2024Secondary}, який охоплює 129 досліджень застосування доповненої реальності у шкільній освіті, засвідчує статистично значущий позитивний вплив AR-технологій на навчальні результати учнів із загальним розміром ефекту g\,=\,0{,}919. Автори встановлюють, що AR-інструменти найбільш ефективні у природничо-математичних дисциплінах, де візуалізація абстрактних понять і моделювання процесів суттєво посилюють практичну спрямованість навчання та підвищують мотивацію до пізнавальної діяльності. Результати мета-аналізу також підтверджують ефективність доповненої реальності на рівні початкової школи, де AR-засоби сприяють підтриманню навчального інтересу молодших школярів та формуванню базових цифрових навичок.

У систематичному огляді Б.~Цзяна зі співавторами~\cite{Nosenko2023Blended}, який охоплює 117 досліджень застосування VR та AR у шкільній STEM-освіті, встановлено, що імерсивні технології суттєво підвищують мотивацію учнів, сприяють візуалізації складних понять і є дієвим інструментом забезпечення безперервності навчання. Автори систематизують підходи до використання віртуальних лабораторій та симуляцій, обґрунтовуючи їх ефективність для навчання природничо-математичних дисциплін, та підкреслюють потенціал імерсивних технологій для організації дистанційного навчання в умовах обмеженого доступу до традиційних освітніх форматів.

Технічний і методичний вимір упровадження імерсивних технологій розкривається у публікації Т.~Цимбалюк і Д.~Федасюка \cite{Tsymbaliuk2024Classification} де здійснено класифікацію освітніх імерсивних середовищ, визначено переваги VR-платформ порівняно з традиційними мультимедійними засобами та окреслено вимоги до апаратного забезпечення. Автори наголошують на необхідності добору рівня занурення з урахуванням вікових і когнітивних особливостей здобувачів освіти.

Вагомий внесок у формування змістового компонента навчання з використанням імерсивних технологій здійснила Ю.~Носенко \cite{Nosenko2025ITE}, яка у своєму дослідженні обґрунтовує підходи до формування змістового компонента навчання в імерсивних освітніх середовищах у закладах загальної середньої освіти. У роботі наголошується на необхідності цілеспрямованого добору та структурування навчального контенту для забезпечення педагогічної ефективності VR/AR-рішень.

Практичні аспекти використання імерсивних технологій у STEM-освіті систематизовано у скопінговому огляді Н.~Пелласа, А.~Денгеля та А.~Крістопулоса \cite{Drokina2025STEM}, який охоплює 41 дослідження та присвячений аналізу застосування імерсивної віртуальної реальності у навчанні природничо-математичних і технічних дисциплін. Автори доводять, що VR-засоби сприяють формуванню дослідницьких навичок, просторового мислення та інтересу до STEM-дисциплін, водночас підкреслюючи необхідність подальшого вдосконалення методичних підходів до їх інтеграції в освітній процес.

Порівняльний аналіз впливу різних типів імерсивних технологій на когнітивне навантаження та мотивацію здійснено у систематичному огляді М.~Пупара зі співавторами~\cite{Poupard2024Systematic}. Автори встановили, що AR-технології знижують когнітивне навантаження та є більш ефективними для навчання початківців порівняно з VR, тоді як VR-середовища підвищують внутрішню мотивацію, але збільшують стороннє когнітивне навантаження. Це має пряме значення для проєктування освітнього процесу підготовки майбутніх учителів, які на початковому етапі опанування імерсивних технологій потребують оптимізації когнітивного навантаження.

Методологічні засади створення імерсивних освітніх ресурсів розглянуто у колективній праці С.~Семерікова зі співавторами \cite{Semerikov2022IERDesign} у якій запропоновано дизайн-модель розроблення VR/AR-засобів навчання з урахуванням дидактичних принципів, когнітивного навантаження та оцінювання ефективності занурювального досвіду. Автори наголошують на важливості педагогічного дизайну й комунікативної взаємодії між учасниками освітнього процесу.

У контексті вищої освіти С.~Кумар та А.~Гораї \cite{Gnedina2025Quality} на основі мета-аналізу досліджень ефективності AR- та VR-втручань встановлюють статистично значущий вплив імерсивних технологій на навчальні результати (d\,=\,0{,}98), підкреслюючи їхній потенціал для підвищення залученості студентів, розвитку критичного мислення та академічної мотивації у закладах вищої освіти. Подібні ідеї простежуються у працях Ю.~Матвієнка \cite{Lytvynova2023Comparative}, де імерсивні технології розглядаються як засіб модернізації освітнього середовища, підвищення інтерактивності навчання та забезпечення доступності освітніх матеріалів через відкриті цифрові платформи.

Стратегічний вимір впливу імерсивних технологій на розвиток креативності здобувачів освіти окреслено у мета-аналізі Дж. Сюна зі співавторами \cite{Kindzhibala2025Creative}, який охоплює 19 досліджень і встановлює загальний розмір ефекту 0{,}82 для впливу імерсивних технологій на креативність студентів. Автори доводять, що VR/AR-засоби виступають дієвим інструментом креативного навчання, спрямованого на формування практикоорієнтованих компетентностей та критичного мислення. Завершує спектр праць другого кластера систематичний огляд Т.~Тене зі співавторами \cite{Medvedieva2024STEAM}, присвячений інтеграції імерсивних технологій у STEM-освіту. Автори визначають VR/AR-технології як каталізатор інтеграції природничих, технічних і гуманітарних знань, що сприяє розвитку інженерного мислення, креативності та міждисциплінарного світогляду учнів.

Дж. Паронг та Р. Е. Майєр \cite{Parong2021Cognitive} аналізують \emph{вплив ІОР на навчальні результати}, а також які \emph{афективні та когнітивні процеси} опосередковують цей вплив.

Праці другого кластера засвідчують поступове зміщення наукового інтересу від концептуального осмислення імерсивних технологій до аналізу реальних освітніх практик їх упровадження на різних рівнях освіти. Дослідження цього напряму демонструють різноманіття педагогічних сценаріїв використання VR/AR-ресурсів -- від початкової й профільної школи до вищої освіти, у межах традиційного, змішаного та дистанційного навчання, а також в умовах соціальних і безпекових викликів. Узагальнені результати дозволяють констатувати, що ефективність імерсивних технологій у навчанні зумовлюється поєднанням дидактичної доцільності, технічної готовності освітнього середовища та методичної підготовленості педагогів. Водночас наявні дослідження зосереджені переважно на вивченні впливу ІОР на здобувачів освіти, тоді як питання підготовки педагогів до \emph{створення} цих ресурсів залишається недостатньо розробленим.

%%% CLUSTER 3 %%%

На цьому підґрунті у \emph{працях третього кластера} фокус наукового аналізу зміщується на \textit{питання професійної педагогічної підготовки та розвитку педагогічного професіоналізму,} у межах яких імерсивні технології розглядаються вже не лише як засіб навчання учнів, а як інструмент формування професійних, цифрових і методичних компетентностей майбутніх учителів. Дослідження цього напряму фокусуються на визначенні педагогічних, організаційних і методичних умов формування готовності майбутніх учителів до проєктування та використання VR/AR/MR-технологій у професійній діяльності, а також на переосмисленні змісту й структури педагогічної освіти в умовах цифрової трансформації.

Питання формування технологічної готовності майбутніх педагогів ґрунтовно розглянуто у систематичному огляді літератури Х.~Мени зі співавторами \cite{Shevchuk2024Readiness}, присвяченому професійній підготовці вчителів засобами доповненої реальності. Автори систематизують підходи до використання AR у педагогічній освіті, встановлюючи, що ця технологія сприяє формуванню як технологічних, так і методичних компетентностей майбутніх учителів. У роботі окреслено дидактичний потенціал AR для створення інтерактивного освітнього середовища та розвитку креативного мислення. Особливу увагу приділено педагогічним умовам підготовки майбутніх учителів, зокрема інтеграції спеціалізованих дисциплін у зміст педагогічної освіти та створенню практикоорієнтованого освітнього середовища.

У контексті модернізації програм педагогічної освіти та розвитку дистанційних форматів підготовки проблему використання імерсивних технологій розглядає Н.~Яремчук \cite{Omelchuk2024Secondary}. Авторка концептуалізує VR/AR/MR-засоби як новітню форму навчальної наочності, що базується на ефекті занурення та суттєво змінює сприймання навчального образу, розширюючи можливості візуалізації освітнього контенту. У дослідженні проаналізовано програмно-технічний потенціал імерсивних технологій, їхні дидактичні можливості, переваги й обмеження у професійній підготовці вчителів початкової школи, а також обґрунтовано організаційно-педагогічні, методичні та психологічні чинники ефективного використання таких інструментів у цифровому освітньому середовищі.

Окрему групу праць третього кластера становлять дослідження, у яких імерсивні технології розглядаються крізь призму галузевої підготовки педагогів. Результати систематичного огляду Дж. Радіанті та співавторів~\cite{Radianti2020VR} засвідчують можливості VR-технологій для моделювання складних процесів, що важко відтворюються в реальних умовах, та підкреслюють значення імерсивних ресурсів для розвитку дослідницьких умінь студентів і підвищення якості підготовки у різних галузях, зокрема природничих науках.

Розширений аналіз можливостей імерсивних технологій у педагогічній освіті здійснено у систематичному огляді Х.~Фернандес-Серо зі співавторами \cite{Malytska2022EU}, які досліджують перспективи розширеної реальності (XR) в освітньому процесі. Автори встановлюють, що XR-технології створюють умови для розвитку критичного мислення, командної взаємодії та креативності, забезпечуючи трансформацію традиційних моделей підготовки фахівців в умовах цифрової економіки.

Роль імерсивних технологій у системі підготовки педагогів обґрунтовано у систематичному огляді М.~Ванг та Л.~Лі \cite{Angelov2025Synergy}, який на матеріалі 52 досліджень аналізує, як VR, AR та MR сприяють удосконаленню педагогічної освіти. Автори встановлюють, що поєднання імерсивних технологій створює умови для адаптивного, інтерактивного й занурювального навчання, акцентуючи увагу на необхідності формування цифрової культури викладачів педагогічних університетів і їхньої готовності до проєктування імерсивних освітніх курсів. У ширшому контексті розвитку екосистеми освітніх технологій імерсивні рішення розглядаються у систематичному огляді Т.~Тене зі співавторами \cite{Medvedieva2024STEAM}, присвяченому інтеграції імерсивних технологій у STEM-освіту. Автори систематизують сучасні підходи до використання VR/AR-засобів як складової EdTech-екосистеми, підкреслюючи їхню значущість для підготовки майбутніх фахівців в умовах трансформації освітніх систем та забезпечення неперервності освіти.

Актуальним напрямом досліджень третього кластера є аналіз компетентнісних розривів у підготовці педагогів до роботи з імерсивними технологіями. Х.~Бельда-Медіна та Х.~Р.~Кальво-Феррер~\cite{BeldaMedina2022TPACK} на основі моделі TPACK встановили, що майбутні вчителі демонструють суттєвий дефіцит на перетині технологічних і педагогічних знань у сфері доповненої реальності, оскільки система педагогічної освіти готує їх переважно як \emph{користувачів} AR-засобів, а не як \emph{розробників}. Дж.~Янг зі співавторами~\cite{Jang2021TAM} підтверджують, що рівень TPACK педагога статистично значуще впливає на сприйняту корисність і намір використовувати AR/VR-технології, що актуалізує потребу цілеспрямованого формування відповідних компетентностей. Перспективний підхід до вирішення цієї проблеми запропоновано С.~Обердорфером зі співавторами~\cite{Oberdorfer2021Interdisciplinary}, які експериментально підтвердили ефективність міждисциплінарної моделі навчання, що поєднує студентів педагогічних та ІТ-спеціальностей у спільному проєктуванні VR/AR-ресурсів, забезпечуючи значне покращення як медіадидактичних, так і технологічних знань обох груп.

Отже, праці третього кластера засвідчують, що імерсивні технології в сучасному науково-педагогічному дискурсі дедалі частіше осмислюються як інструмент трансформації системи професійної педагогічної підготовки. Водночас виразно простежується ключова суперечність: наявні дослідження здебільшого зосереджені на підготовці вчителів до \emph{використання} VR/AR/MR-засобів, тоді як проблема формування їхньої готовності до \emph{розробки} імерсивних освітніх ресурсів залишається недостатньо розробленою~\cite{BeldaMedina2022TPACK, Oberdorfer2021Interdisciplinary}. На цьому тлі у працях четвертого кластера фокус наукового аналізу спрямовується на проблематику використання імерсивних технологій в інклюзивній освіті, де питання доступності, адаптивності та безбар'єрності освітнього середовища набувають визначального значення.

%%% CLUSTER 4 %%%

У межах четвертого кластера наукових джерел у центрі уваги перебуває проблематика використання імерсивних освітніх ресурсів в умовах інклюзивної парадигми освіти. Дослідження цього напряму зосереджені на осмисленні VR/AR-технологій як інструментів забезпечення доступності навчального контенту, компенсації сенсорних, когнітивних і комунікативних обмежень, а також створення безпечного, адаптивного й варіативного освітнього середовища для здобувачів освіти з різними освітніми потребами. У цьому контексті імерсивні технології розглядаються не лише як дидактичне нововведення, а як механізм реалізації принципів безбар'єрності, універсального дизайну навчання та рівного доступу до якісної освіти.

Можливості застосування імерсивних технологій у системі інклюзивної освіти ґрунтовно розкрито у систематичному огляді Х.~Халкіадакіса зі співавторами~\cite{Chyhryna2023Inclusive}, присвяченому аналізу впливу штучного інтелекту та віртуальної реальності на забезпечення інклюзивності в освіті. Автори обґрунтовують ефективність VR-інструментів у роботі з учнями з особливими освітніми потребами, встановлюючи, що імерсивні середовища у поєднанні з адаптивними AI-алгоритмами дозволяють мінімізувати комунікативні бар'єри, знизити рівень тривожності, сприяти соціальній інтеграції та створити безпечний навчальний простір, адаптований до індивідуальних можливостей здобувачів освіти. Використання VR-платформ у поєднанні із засобами штучного інтелекту розширює можливості включення дітей з особливими освітніми потребами у спільну навчальну діяльність, сприяє розвитку комунікативних навичок і забезпечує адаптацію освітнього процесу до індивідуальних освітніх траєкторій.

Комплексний аналіз застосування імерсивних технологій в інклюзивній цифровій освіті здійснено у систематичному огляді К.~Поджіанті зі співавторами~\cite{Poggianti2025Inclusive}, які встановили, що лише 40{,}6\,\% досліджень у цій сфері явно застосовують рамку універсального дизайну навчання (UDL), що свідчить про недостатню педагогічну обґрунтованість значної частини розробок і підкреслює потребу у цілеспрямованій підготовці розробників ІОР до врахування принципів інклюзивного дизайну.

У новітніх дослідженнях українських науковців інклюзивний потенціал імерсивних технологій дедалі частіше розглядається у взаємозв'язку з іншими провідними освітніми напрямами, зокрема STEM-освітою, гуманітарною підготовкою та емоційно-ціннісним розвитком. Такий міждисциплінарний підхід засвідчує розширення наукового бачення імерсивних технологій як комплексного ресурсу, здатного одночасно підтримувати пізнавальний, соціальний і особистісний розвиток здобувачів освіти в умовах інклюзивного освітнього середовища.

Узагальнення наукових позицій дозволяє констатувати, що VR/AR-ресурси в умовах інклюзивної освіти виконують не лише компенсаторну функцію, а й виступають засобом розширення освітніх можливостей, соціальної взаємодії та особистісного розвитку здобувачів освіти. Вони забезпечують реалізацію принципів безбар'єрності, універсального дизайну навчання й рівного доступу до якісної освіти, що істотно посилює гуманістичний потенціал сучасного освітнього середовища. На цьому підґрунті у працях п'ятого кластера науковий аналіз зосереджується на галузевих, порівняльних і узагальнювальних підходах до використання імерсивних технологій в освіті, що дає змогу співвіднести національний і міжнародний досвід, а також виявити спільні тенденції та специфічні особливості розвитку цього напряму.

%%% CLUSTER 5 %%%

\emph{У межах п'ятого кластера наукових} джерел зосереджено галузеві, порівняльні та узагальнювальні підходи до використання імерсивних технологій в освіті, що дозволяють простежити еволюцію цього напряму в міжнародному та вітчизняному контекстах, а також виявити міждисциплінарні тенденції його розвитку. Дослідження цього кластера репрезентують імерсивні технології як стратегічний інструмент цифрової трансформації освіти, здатний інтегрувати педагогічні, технологічні й гуманітарні виміри навчання.

Значне місце в межах кластера посідають праці, присвячені узагальненню міжнародного досвіду використання імерсивних технологій. Так, у систематичному огляді Х.~Фернандес-Серо зі співавторами \cite{Malytska2022EU} проаналізовано можливості розширеної реальності (XR) в освітньому процесі, зокрема у контексті підготовки педагогів, розвитку цифрової грамотності та STEM-освіти. Автори підкреслюють, що результативність упровадження імерсивних рішень забезпечується узгодженістю технологічних можливостей XR-платформ із дидактичними цілями та готовністю освітніх систем до їх інтеграції.

Порівняльний аналіз ефективності імерсивного навчання здійснено у мета-аналізі М.~Кобана зі співавторами \cite{Lytvynova2023Comparative}, який охоплює 48 досліджень і присвячений потенціалу імерсивної віртуальної реальності для підвищення якості навчання. Автори встановлюють статистично значущий позитивний вплив VR на навчальні результати, визначаючи імерсивні технології як стратегічний напрям цифровізації освіти, що потребує поєднання інфраструктурного забезпечення та системної підготовки педагогічних кадрів. Подібні ідеї розвивають Х.~Мена зі співавторами \cite{Shevchuk2024Readiness} у систематичному огляді літератури, присвяченому професійній підготовці вчителів засобами доповненої реальності. Автори узагальнюють міжнародний досвід використання AR у педагогічній освіті, підкреслюючи, що у розвинених освітніх системах ці інструменти вже інтегровані до стандартів педагогічної майстерності, водночас зберігаючи значний потенціал для подальшого розширення у системах, де вони перебувають на етапі становлення.

Окрему групу праць п'ятого кластера становлять узагальнювальні та аналітичні видання, спрямовані на систематизацію наукового доробку та освітніх практик. У мета-аналізі Ч.~Чанга зі співавторами \cite{Lytvynova2023Health}, який охоплює 134 дослідження за десять років, здійснено комплексний огляд застосування доповненої реальності в освіті, що засвідчує стійку тенденцію до інтеграції імерсивних технологій у різні рівні та форми навчання. У мета-аналізі М.~Кобана зі співавторами \cite{Lytvynova2023Comparative}, що охоплює 48 досліджень потенціалу імерсивної віртуальної реальності, обґрунтовано можливості створення інтерактивних освітніх просторів, наголошуючи на важливості міждисциплінарної співпраці педагогів, методистів і розробників цифрових платформ. У колективній монографії С.~Литвинової, Н.~Сороко та С.~Баценко \cite{Lytvynova2023Model} запропоновано модель проєктування освітнього середовища з використанням VR/AR-засобів, ефективність якої пов'язується з дотриманням принципів активності, безпечності та динамічної педагогічної взаємодії.

Важливим складником п'ятого кластера є дослідження галузевого та міждисциплінарного застосування імерсивних технологій. Так, у систематичному огляді А.~Рамтогула та К.~Кхедо \cite{Granchak2023Heritage} проаналізовано застосування систем доповненої реальності у сфері збереження та популяризації культурної спадщини, підкреслюючи потенціал AR-технологій для створення віртуальних музеїв і 3D-експозицій, що розширюють доступ до культурних об'єктів. У мета-аналізі Ц.~Чень зі співавторами \cite{Boiko2024Immersive}, який охоплює 17 досліджень, доведено ефективність технологій розширеної реальності у навчанні іноземних мов із загальним розміром ефекту g\,=\,0{,}825, що свідчить про значний позитивний вплив XR-середовищ на розвиток мовних навичок, зокрема письмової комунікації, занурення у автентичне мовленнєве середовище та подолання комунікативних бар'єрів. Застосування VR-технологій для розвитку дизайнерського мислення аналізується у літературному дослідженні М.~Лю зі співавторами \cite{Gordeev2025Printing}, які систематизують підходи до формування навичок дизайн-мислення у віртуальних середовищах та наголошують на потенціалі VR для розвитку нових професійних компетентностей.

Розширення гуманітарного виміру імерсивних технологій простежується у працях С.~Доценка і В.~Чжен \cite{Dotsenko2023Immersive} де описано феномен «імерсивного шоу» як синтез мистецтва, театру та цифрових технологій, а також у дослідженнях мовної освіти, де Н.~Сороко та О.~Гаєвська \cite{Soroko2021Eastern} обґрунтовують ефективність VR-модулів у викладанні східних мов. Питання змістового дизайну уроку з використанням імерсивних технологій в умовах змішаного навчання розглядають С.~Баценко та В.~Горбаченко \cite{Batsenko2025Blended}, обґрунтовуючи підходи до структурування навчального заняття із залученням VR/AR-засобів і наголошуючи на необхідності узгодження технологічних рішень із дидактичними цілями.

У комплексному огляді А.~Аль-Ансі зі співавторами \cite{Sabodashko2024Humanities}, який охоплює аналіз 1536 публікацій, проаналізовано сучасний розвиток AR та VR в освіті, зокрема в гуманітарній галузі. Автори наголошують, що VR/AR-засоби не замінюють традиційне навчання, але виступають його важливим доповненням, особливо у процесі формування медіа- й цифрової грамотності, а також критичного осмислення навчального контенту. Галузевий вимір імерсивного навчання у технічній освіті розкрито у дослідженні Л.~Нікітіної та Н.~Дженіюк \cite{Nikitina2023Telecom} де імерсивні технології розглядаються як ефективний інструмент підготовки фахівців у сфері телекомунікацій. Авторки доводять, що інтерактивні 3D-моделі та графічні симуляції сприяють глибшому засвоєнню теоретичних положень дисциплін, розвитку просторового мислення та формуванню навичок технічного аналізу.

Цілісну картину розвитку напряму доповнює бібліометричне дослідження С.~Гарнала зі співавторами~\cite{Harnal2023Bibliometric}, які на основі аналізу 3909 публікацій здійснили тематичне картування досліджень доповненої реальності в освіті, виявивши чотири основні тематичні кластери та встановивши стійку тенденцію до зростання міждисциплінарності досліджень. Завершує огляд систематичний огляд А.~Рамтогула та К.~Кхедо \cite{Granchak2023Heritage}, у якому AR-технології розглядаються як засіб збереження та відтворення культурної спадщини, здатний посилювати доступність культурних і освітніх об'єктів та активну участь здобувачів освіти у пізнавальному процесі.

Загалом праці п'ятого кластера репрезентують імерсивні технології як міждисциплінарний і стратегічний ресурс розвитку сучасної освіти, що виходить за межі окремих рівнів і галузей навчання. Узагальнення галузевого, порівняльного та аналітичного досвіду засвідчує, що ефективність упровадження VR/AR-рішень зумовлюється не лише технічними інноваціями, а й системністю освітньої політики, підготовкою педагогічних кадрів, міждисциплінарною взаємодією та орієнтацією на гуманістичні цінності освіти.

%%% Figure 1.2 %%%

З метою систематизації наукових підходів до вивчення імерсивних освітніх ресурсів та виявлення логіки їх проблемно-тематичного осмислення в сучасному науково-педагогічному дискурсі доцільно представити деталізовану модель кластерного групування досліджень, що конкретизує рис.~\ref{fig:1.2}.

\begin{figure}[!h]
\centering
%\resizebox{\textwidth}{!}{%
\begin{tikzpicture}

% Title
\node[dia/header=14cm] (title) {ПРОБЛЕМНО-ТЕМАТИЧНІ КЛАСТЕРИ};

% Upper-left cluster
\node[dia/lrbox=6cm, below=1.5cm of title.south, xshift=-4cm] (c1) {%
  \textbf{Загальнотеоретичний і дидактичний кластер}\\[3pt]
  \textbullet~осмислення ІОР;\\
  \textbullet~трансформація змісту, форм і методів навчання;\\
  \textbullet~дидактичний потенціал ІТ%
};

% Upper-right cluster
\node[dia/lrbox=6cm, below=1.5cm of title.south, xshift=4cm] (c2) {%
  \textbf{Освітні практики}\\[3pt]
  \textbullet~різні ланки освіти;\\
  \textbullet~педагогічні умови;\\
  \textbullet~організаційні моделі;\\
  \textbullet~дидактичні сценарії%
};

% Lower-left cluster
\node[dia/lrbox=6cm, below=1.2cm of c1] (c3) {%
  \textbf{Галузеві, порівняльні та узагальнювальні підходи}\\[3pt]
  \textbullet~міждисциплінарність;\\
  \textbullet~порівняльний аналіз міжнародного і вітчизняного досвіду;\\
  \textbullet~узагальнення наукових та прикладних напрацювань%
};

% Lower-right cluster
\node[dia/lrbox=6cm, below=1.2cm of c2] (c4) {%
  \textbf{Професійна педагогічна підготовка}\\[3pt]
  \textbullet~формування професійних і цифрових компетентностей;\\
  \textbullet~проєктування освітніх рішень;\\
  \textbullet~розвиток професіоналізму%
};

% Tree structure: title → row 1, row 1 → row 2
\draw[dia/line] (title.south) -- (c1.north);
\draw[dia/line] (title.south) -- (c2.north);
\draw[dia/line] (c1.south) -- (c3.north);
\draw[dia/line] (c2.south) -- (c4.north);

% Convergence point
\coordinate (conv) at ($(c3.south)!0.5!(c4.south) + (0,-1.5cm)$);
\node[dia/label] at (conv) (convlabel) {Точка перетину};

% Lines from row 2 boxes to convergence (direct, no overlap)
\draw[dia/line] (c3.south) -- (convlabel.north);
\draw[dia/line] (c4.south) -- (convlabel.north);

% Lines from row 1 boxes to convergence (routed outside row 2)
\draw[dia/line] (c1.south west) -- ([xshift=-0.4cm]c3.south west) -- (convlabel.north west);
\draw[dia/line] (c2.south east) -- ([xshift=0.4cm]c4.south east) -- (convlabel.north east);

% Bottom cluster
\node[dia/lrbox=6cm, below=1.2cm of convlabel] (inclusive) {%
  \textbf{Інклюзивний освітній вимір}\\[3pt]
  \textbullet~доступність;\\
  \textbullet~підтримка освітніх потреб;\\
  \textbullet~безбар'єрність%
};

% Arrow from convergence to bottom box
\draw[dia/arrow] (convlabel.south) -- (inclusive.north);

\end{tikzpicture}%
%}

\caption{Змістове наповнення проблемно-тематичних кластерів науково-педагогічного осмислення імерсивних технологій в освіті}
\label{fig:1.2}
\end{figure}

%%% PART D: Integrative Synthesis %%%

Подана модель відображає взаємопов'язаність проблемно-тематичних кластерів імерсивних технологій в освіті та демонструє інклюзивний освітній вимір як наскрізну точку перетину дидактичних, практичних, професійно-педагогічних і галузево-порівняльних підходів.

Проведений аналіз п'яти проблемно-тематичних кластерів засвідчує, що наукове осмислення імерсивних технологій в освіті пройшло еволюцію від концептуалізації (кластер~1) через практичне впровадження (кластер~2) до питань педагогічної підготовки (кластер~3), інклюзивного застосування (кластер~4) та міждисциплінарних узагальнень (кластер~5). Водночас аналіз виявляє фундаментальну суперечність: за наявності переконливих доказів ефективності ІОР (розміри ефекту g\,=\,0{,}82--0{,}98 за результатами мета-аналізів, табл.~\ref{tab:meta_analysis}) та розвинутої теоретичної бази (CAMIL~\cite{Makransky2021CAMIL}, теорія когнітивного навантаження~\cite{Buchner2021CogLoad, Poupard2024Systematic}) система педагогічної освіти залишається зорієнтованою на підготовку вчителів як \emph{користувачів} імерсивних технологій, а не їх \emph{творців}~\cite{BeldaMedina2022TPACK}.

Ця суперечність зумовлює потребу у формуванні цілісного компетентнісного профілю вчителя інформатики як розробника ІОР (рис.~\ref{fig:competency_profile}), що поєднує технологічну компетентність (WebGL, Three.js, MindAR, TensorFlow.js, WebXR), педагогічну компетентність (дидактичний дизайн, сценарування занурення, інклюзивний дизайн) та інтегративну компетентність (проєктування ІОР, міждисциплінарний підхід, рефлексивна практика).

\begin{figure}[!h]
\centering
%\resizebox{\textwidth}{!}{%
\begin{tikzpicture}

% Title
\node[dia/header=14cm] (title) {КОМПЕТЕНТНІСНИЙ ПРОФІЛЬ\\ УЧИТЕЛЯ ІНФОРМАТИКИ ЯК РОЗРОБНИКА ІОР};

% Three competency domain boxes
\node[dia/lrbox=4cm, below=1.5cm of title.south, xshift=-5cm] (tech) {%
  \textbf{Технологічна компетентність}\\[3pt]
  \textbullet~WebGL / Three.js;\\
  \textbullet~MindAR / AR.js;\\
  \textbullet~TensorFlow.js;\\
  \textbullet~WebXR Device API;\\
  \textbullet~3D-моделювання%
};

\node[dia/lrbox=4cm, below=1.5cm of title.south] (ped) {%
  \textbf{Педагогічна \\компетентність}\\[3pt]
  \textbullet~дидактичний дизайн ІОР;\\
  \textbullet~сценарування занурення;\\
  \textbullet~оцінювання досвіду;\\
  \textbullet~інклюзивний дизайн;\\
  \textbullet~медіаграмотність%
};

\node[dia/lrbox=4cm, below=1.5cm of title.south, xshift=5cm] (integ) {%
  \textbf{Інтегративна компетентність}\\[3pt]
  \textbullet~проєктування ІОР;\\
  \textbullet~міждисциплінар\-ність;\\
  \textbullet~рефлексивна практика;\\
  \textbullet~громадянська освіта;\\
  \textbullet~командна взаємодія%
};

% Arrows from title to domains
\draw[dia/arrow] (title.south) -- (tech.north);
\draw[dia/arrow] (title.south) -- (ped.north);
\draw[dia/arrow] (title.south) -- (integ.north);

% Convergence box below
\node[dia/rbox=10cm, below=1.5cm of ped.south] (result) {%
  Готовність до розробки імерсивних освітніх ресурсів%
};

% Arrows from domains to result
\draw[dia/arrow] (tech.south) -- (result.north);
\draw[dia/arrow] (ped.south) -- (result.north);
\draw[dia/arrow] (integ.south) -- (result.north);

% Bidirectional arrows between domains
\draw[dia/darrow] (tech.east) -- (ped.west);
\draw[dia/darrow] (ped.east) -- (integ.west);

\end{tikzpicture}%
%}

\caption{Компетентнісний профіль учителя інформатики як розробника імерсивних освітніх ресурсів}
\label{fig:competency_profile}
\end{figure}

Реалізація зазначеного компетентнісного профілю потребує опори на відповідний технологічний стек WebAR (рис.~\ref{fig:webar_stack}), який утворює багаторівневу архітектуру від браузерної платформи (WebGL, WebRTC) через JavaScript-бібліотеки (Three.js, TensorFlow.js, WebXR Device API) та WebAR-рушії (MindAR, AR.js) до рівня освітнього застосування~-- створення імерсивних освітніх ресурсів.

\begin{figure}[!h]
\centering
%\resizebox{\textwidth}{!}{%
\begin{tikzpicture}

% Layer 1 (bottom): Browser platform
\node[dia/lrbox=13cm] (browser) {%
  \textbf{Браузерна платформа}\\[3pt]
  \textbullet~WebGL --- рендеринг 3D-графіки;\\
  \textbullet~WebRTC --- доступ до камери та мікрофона;\\
  \textbullet~Canvas API --- 2D-графіка та композиція%
};

% Layer 2: Libraries
\node[dia/lrbox=13cm, above=1cm of browser] (libs) {%
  \textbf{JavaScript-бібліотеки}\\[3pt]
  \textbullet~Three.js / A-Frame --- 3D-візуалізація та сцени;\\
  \textbullet~TensorFlow.js / Teachable Machine --- машинне навчання;\\
  \textbullet~WebXR Device API --- доступ до XR-пристроїв%
};

% Layer 3: AR engines
\node[dia/lrbox=13cm, above=1cm of libs] (engines) {%
  \textbf{WebAR-рушії та фреймворки}\\[3pt]
  \textbullet~MindAR --- відстеження зображень, облич;\\
  \textbullet~AR.js --- маркерне та GPS-орієнтоване AR;\\
  \textbullet~8th Wall / Zappar --- комерційні рішення%
};

% Layer 4 (top): Educational application
\node[dia/rbox=13cm, above=1cm of engines] (edu) {%
  \textbf{Освітнє застосування}: імерсивні освітні ресурси (ІОР)%
};

% Arrows between layers
\draw[dia/arrow] (browser.north) -- (libs.south);
\draw[dia/arrow] (libs.north) -- (engines.south);
\draw[dia/arrow] (engines.north) -- (edu.south);

\end{tikzpicture}%
%}

\caption{Технологічний стек WebAR для розробки імерсивних освітніх ресурсів}
\label{fig:webar_stack}
\end{figure}

Спроби вирішення означеної суперечності представлено, зокрема, у працях С.~О.~Семерікова зі співавторами, де запропоновано дизайн-модель розроблення імерсивних освітніх ресурсів~\cite{Semerikov2022IERDesign} та розроблено методику навчання WebAR з інтеграцією машинного навчання~\cite{Semerikov2024WebAR_EduDim}. Однак цілісна методика підготовки майбутніх учителів інформатики до розробки ІОР, що враховувала б теоретичні засади CAMIL, принципи управління когнітивним навантаженням та модель TPACK, залишається нерозробленою, що зумовлює необхідність подальшого дослідження сутнісних характеристик імерсивних освітніх ресурсів (підрозділ~\ref{subsec:1.2}).
